\section*{Exercise 4}
\textit{Deletion of Unused Code, Meaningful Comments, and JavaDocs}

\noindent As the final task of this assignment, we were asked to clean up the codebase and ensure it follows good software engineering practices, including proper documentation and modular structure.

\noindent In addition to restructuring the project into packages and adding Javadoc comments, we also removed unused or redundant code and added explanatory comments where necessary. Rather than repeating what was done in detail, it is more useful at this point to summarize the main benefits of adhering to these guidelines:

\begin{itemize}
\item \textbf{Clean formatting:} Improves code readability and helps developers navigate and understand the structure more quickly.

\item \textbf{Non-redundant code:} Removing unnecessary or duplicate logic reduces complexity and minimizes potential sources of error.

\item \textbf{Comments for non-trivial functions:} Improve comprehensibility, especially for algorithms or logic that are not immediately self-explanatory.

\item \textbf{Java documentation (JavaDocs):} Provides standardized, accessible documentation for all public classes and methods, supporting better collaboration and easier future maintenance.

\item \textbf{Consistent package structure:} Promotes modularity and separates concerns, which simplifies testing, debugging, and future extensions.
\end{itemize}

\noindent Altogether, this clean-up step ensures that the code remains maintainable, extendable, and easier to work with for both current and future developers.
